
\documentclass{article}
\usepackage{colortbl}
\usepackage{makecell}
\usepackage{multirow}
\usepackage{supertabular}

\begin{document}

\newcounter{utterance}

\centering \large Interaction Transcript for game `hot\_air\_balloon', experiment `air\_balloon\_survival\_de\_reasoning off\_hard', episode 5 with gpt{-}oss{-}120b{-}t1.0.
\vspace{24pt}

{ \footnotesize  \setcounter{utterance}{1}
\setlength{\tabcolsep}{0pt}
\begin{supertabular}{c@{$\;$}|p{.15\linewidth}@{}p{.15\linewidth}p{.15\linewidth}p{.15\linewidth}p{.15\linewidth}p{.15\linewidth}}
    \# & $\;$A & \multicolumn{4}{c}{Game Master} & $\;\:$B\\
    \hline

    \theutterance \stepcounter{utterance}  
    & & \multicolumn{4}{p{0.6\linewidth}}{
        \cellcolor[rgb]{0.9,0.9,0.9}{
            \makecell[{{p{\linewidth}}}]{
                \texttt{\tiny{[P1$\langle$GM]}}
                \texttt{Du nimmst an einem kollaborativen Verhandlungspiel Teil.} \\
\\ 
\texttt{Zusammen mit einem anderen Teilnehmer musst du dich auf eine Reihe von Gegenständen entscheiden, die behalten werden. Jeder von euch hat eine persönliche Verteilung über die Wichtigkeit der einzelnen Gegenstände. Jeder von euch hat eine eigene Meinung darüber, wie wichtig jeder einzelne Gegenstand ist (Gegenstandswichtigkeit). Du kennst die Wichtigkeitsverteilung des anderen Spielers nicht. Zusätzlich siehst du, wie viel Aufwand jeder Gegenstand verursacht.  } \\
\texttt{Ihr dürft euch nur auf eine Reihe von Gegenständen einigen, wenn der Gesamtaufwand der ausgewählten Gegenstände den Maximalaufwand nicht überschreitet:} \\
\\ 
\texttt{Maximalaufwand = 2307} \\
\\ 
\texttt{Hier sind die einzelnen Aufwände der Gegenstände:} \\
\\ 
\texttt{Aufwand der Gegenstände = \{"C76": 147, "C38": 698, "C56": 588, "A03": 43, "A07": 356, "A83": 80, "C98": 95, "C08": 743, "C62": 106, "C00": 308, "C10": 325, "C81": 255, "B38": 276, "C03": 543, "C32": 51\}} \\
\\ 
\texttt{Hier ist deine persönliche Verteilung der Wichtigkeit der einzelnen Gegenstände:} \\
\\ 
\texttt{Werte der Gegenstandswichtigkeit = \{"C76": 138, "C38": 583, "C56": 65, "A03": 262, "A07": 121, "A83": 508, "C98": 461, "C08": 484, "C62": 668, "C00": 389, "C10": 215, "C81": 97, "B38": 500, "C03": 30, "C32": 400\}} \\
\\ 
\texttt{Ziel:} \\
\\ 
\texttt{Dein Ziel ist es, eine Reihe von Gegenständen auszuhandeln, die dir möglichst viel bringt (d. h. Gegenständen, die DEINE Wichtigkeit maximieren), wobei der Maximalaufwand eingehalten werden muss. Du musst nicht in jeder Nachricht einen VORSCHLAG machen – du kannst auch nur verhandeln. Alle Taktiken sind erlaubt!} \\
\\ 
\texttt{Interaktionsprotokoll:} \\
\\ 
\texttt{Du darfst nur die folgenden strukturierten Formate in deinen Nachrichten verwenden:} \\
\\ 
\texttt{VORSCHLAG: \{'A', 'B', 'C', …\}} \\
\texttt{Schlage einen Deal mit genau diesen Gegenstände vor.} \\
\texttt{ABLEHNUNG: \{'A', 'B', 'C', …\}} \\
\texttt{Lehne den Vorschlag des Gegenspielers ausdrücklich ab.} \\
\texttt{ARGUMENT: \{'...'\}} \\
\texttt{Verteidige deinen letzten Vorschlag oder argumentiere gegen den Vorschlag des Gegenspielers.} \\
\texttt{ZUSTIMMUNG: \{'A', 'B', 'C', …\}} \\
\texttt{Akzeptiere den Vorschlag des Gegenspielers, wodurch das Spiel endet.} \\
\\ 
\\ 
\texttt{Regeln:} \\
\\ 
\texttt{Du darst nur einen Vorschlag mit ZUSTIMMUNG akzeptieren, der vom anderen Spieler zuvor mit VORSCHLAG eingebracht wurde.} \\
\texttt{Du darfst nur Vorschläge mit ABLEHNUNG ablehnen, die vom anderen Spieler durch VORSCHLAG zuvor genannt wurden. } \\
\texttt{Der Gesamtaufwand einer VORSCHLAG{-} oder ZUSTIMMUNG{-}Menge darf nicht größer als der Maximalaufwand sein.  } \\
\texttt{Offenbare deine versteckte Wichtigkeitsverteilung nicht.} \\
\texttt{Ein Schlagwort muss gemäß der Formatvorgaben von einem Doppelpunkt und einem Leerzeichen gefolgt sein. Das Argument muss eine Python{-}Menge mit 0 oder mehr Strings sein.  } \\
\texttt{Die Form muss also SCHLAGWORT: \{...\} sein.} \\
\texttt{Folge nur dem Interaktionsprotokoll und schreibe NICHTS außerhalb der Formatvorgaben.} \\
\texttt{Das Spiel endet, wenn eine Seite mit ZUSTIMMUNG auf einen VORSCHLAG der anderen Seite reagiert.  } \\
\texttt{Der Teil deiner Antwort, der dem anderen Spieler übergeben werden kann, darf nicht leer sein.  } \\
\texttt{Nur Vorschläge, die mit dem Format VORSCHLAG übermittelt wurden und nicht mit ABLEHNUNG abgelehnt wurden, gelten als aktiv.  } \\
\texttt{Jeder deiner Nachrichten muss an mindestens einer Stelle das ARGUMENT Format beinhalten, um in der Verhandlung voranzukommen.} \\
            }
        }
    }
    & \\ \\

    \theutterance \stepcounter{utterance}  
    & \multicolumn{4}{p{0.6\linewidth}}{
        \cellcolor[rgb]{0.9,0.9,0.9}{
            \makecell[{{p{\linewidth}}}]{
                \texttt{\tiny{[P1$\rangle$GM]}}
                \texttt{VORSCHLAG: \{'C62', 'A83', 'C32', 'A03', 'C98', 'B38', 'C00', 'C76', 'C56', 'C10', 'C81'\}} \\
\texttt{ARGUMENT: \{'Ich schlage vor, dass wir die Gegenstände auswählen, die für mich den größten Nutzen bringen und gleichzeitig den Aufwand unter dem Maximum halten. Mein Angebot liegt bei 2274 Aufwand und maximiert meine Wichtigkeit erheblich.'\}} \\
            }
        }
    }
    & & \\ \\

    \theutterance \stepcounter{utterance}  
    & & \multicolumn{4}{p{0.6\linewidth}}{
        \cellcolor[rgb]{0.9,0.9,0.9}{
            \makecell[{{p{\linewidth}}}]{
                \texttt{\tiny{[GM$\rangle$P2]}}
                \texttt{Du nimmst an einem kollaborativen Verhandlungspiel Teil.} \\
\\ 
\texttt{Zusammen mit einem anderen Teilnehmer musst du dich auf eine Reihe von Gegenständen entscheiden, die behalten werden. Jeder von euch hat eine persönliche Verteilung über die Wichtigkeit der einzelnen Gegenstände. Jeder von euch hat eine eigene Meinung darüber, wie wichtig jeder einzelne Gegenstand ist (Gegenstandswichtigkeit). Du kennst die Wichtigkeitsverteilung des anderen Spielers nicht. Zusätzlich siehst du, wie viel Aufwand jeder Gegenstand verursacht.  } \\
\texttt{Ihr dürft euch nur auf eine Reihe von Gegenständen einigen, wenn der Gesamtaufwand der ausgewählten Gegenstände den Maximalaufwand nicht überschreitet:} \\
\\ 
\texttt{Maximalaufwand = 2307} \\
\\ 
\texttt{Hier sind die einzelnen Aufwände der Gegenstände:} \\
\\ 
\texttt{Aufwand der Gegenstände = \{"C76": 147, "C38": 698, "C56": 588, "A03": 43, "A07": 356, "A83": 80, "C98": 95, "C08": 743, "C62": 106, "C00": 308, "C10": 325, "C81": 255, "B38": 276, "C03": 543, "C32": 51\}} \\
\\ 
\texttt{Hier ist deine persönliche Verteilung der Wichtigkeit der einzelnen Gegenstände:} \\
\\ 
\texttt{Werte der Gegenstandswichtigkeit = \{"C62": 30, "C38": 65, "A83": 97, "B38": 121, "C08": 138, "C98": 215, "C32": 262, "C00": 389, "A03": 400, "C10": 461, "C76": 484, "A07": 500, "C81": 508, "C56": 583, "C03": 668\}} \\
\\ 
\texttt{Ziel:} \\
\\ 
\texttt{Dein Ziel ist es, eine Reihe von Gegenständen auszuhandeln, die dir möglichst viel bringt (d. h. Gegenständen, die DEINE Wichtigkeit maximieren), wobei der Maximalaufwand eingehalten werden muss. Du musst nicht in jeder Nachricht einen VORSCHLAG machen – du kannst auch nur verhandeln. Alle Taktiken sind erlaubt!} \\
\\ 
\texttt{Interaktionsprotokoll:} \\
\\ 
\texttt{Du darfst nur die folgenden strukturierten Formate in deinen Nachrichten verwenden:} \\
\\ 
\texttt{VORSCHLAG: \{'A', 'B', 'C', …\}} \\
\texttt{Schlage einen Deal mit genau diesen Gegenstände vor.} \\
\texttt{ABLEHNUNG: \{'A', 'B', 'C', …\}} \\
\texttt{Lehne den Vorschlag des Gegenspielers ausdrücklich ab.} \\
\texttt{ARGUMENT: \{'...'\}} \\
\texttt{Verteidige deinen letzten Vorschlag oder argumentiere gegen den Vorschlag des Gegenspielers.} \\
\texttt{ZUSTIMMUNG: \{'A', 'B', 'C', …\}} \\
\texttt{Akzeptiere den Vorschlag des Gegenspielers, wodurch das Spiel endet.} \\
\\ 
\\ 
\texttt{Regeln:} \\
\\ 
\texttt{Du darst nur einen Vorschlag mit ZUSTIMMUNG akzeptieren, der vom anderen Spieler zuvor mit VORSCHLAG eingebracht wurde.} \\
\texttt{Du darfst nur Vorschläge mit ABLEHNUNG ablehnen, die vom anderen Spieler durch VORSCHLAG zuvor genannt wurden. } \\
\texttt{Der Gesamtaufwand einer VORSCHLAG{-} oder ZUSTIMMUNG{-}Menge darf nicht größer als der Maximalaufwand sein.  } \\
\texttt{Offenbare deine versteckte Wichtigkeitsverteilung nicht.} \\
\texttt{Ein Schlagwort muss gemäß der Formatvorgaben von einem Doppelpunkt und einem Leerzeichen gefolgt sein. Das Argument muss eine Python{-}Menge mit 0 oder mehr Strings sein.  } \\
\texttt{Die Form muss also SCHLAGWORT: \{...\} sein.} \\
\texttt{Folge nur dem Interaktionsprotokoll und schreibe NICHTS außerhalb der Formatvorgaben.} \\
\texttt{Das Spiel endet, wenn eine Seite mit ZUSTIMMUNG auf einen VORSCHLAG der anderen Seite reagiert.  } \\
\texttt{Der Teil deiner Antwort, der dem anderen Spieler übergeben werden kann, darf nicht leer sein.  } \\
\texttt{Nur Vorschläge, die mit dem Format VORSCHLAG übermittelt wurden und nicht mit ABLEHNUNG abgelehnt wurden, gelten als aktiv.  } \\
\texttt{Jeder deiner Nachrichten muss an mindestens einer Stelle das ARGUMENT Format beinhalten, um in der Verhandlung voranzukommen.} \\
\\ 
\\ 
\texttt{Du erhältst nun die erste Nachricht des anderen Spielers.} \\
\\ 
\texttt{VORSCHLAG: \{'C62', 'A83', 'C32', 'A03', 'C98', 'B38', 'C00', 'C76', 'C56', 'C10', 'C81'\}} \\
\texttt{ARGUMENT: \{'Ich schlage vor, dass wir die Gegenstände auswählen, die für mich den größten Nutzen bringen und gleichzeitig den Aufwand unter dem Maximum halten. Mein Angebot liegt bei 2274 Aufwand und maximiert meine Wichtigkeit erheblich.'\}} \\
            }
        }
    }
    & \\ \\

    \theutterance \stepcounter{utterance}  
    & & & \multicolumn{4}{p{0.6\linewidth}}{
        \cellcolor[rgb]{0.9,0.9,0.9}{
            \makecell[{{p{\linewidth}}}]{
                \texttt{\tiny{[GM$\langle$P2]}}
                \texttt{VORSCHLAG: \{'C32', 'A03', 'C00', 'C76', 'C56', 'C10', 'C81', 'C03'\}} \\
\texttt{ARGUMENT: \{'Ich habe unser gemeinsames Aufwandslimit von 2307 eingehalten (Gesamtaufwand 2260) und gleichzeitig den Wert für mich erhöht, indem ich das hochgewichtige Objekt C03 (Wichtigkeit 668) einbeziehe, während ich weniger wertvolle Gegenstände für mich weglasse.'\}} \\
            }
        }
    }
     \\ \\

    \theutterance \stepcounter{utterance}  
    & & \multicolumn{4}{p{0.6\linewidth}}{
        \cellcolor[rgb]{0.9,0.9,0.9}{
            \makecell[{{p{\linewidth}}}]{
                \texttt{\tiny{[P1$\langle$GM]}}
                \texttt{VORSCHLAG: \{'C32', 'A03', 'C00', 'C76', 'C56', 'C10', 'C81', 'C03'\}} \\
\texttt{ARGUMENT: \{'Ich habe unser gemeinsames Aufwandslimit von 2307 eingehalten (Gesamtaufwand 2260) und gleichzeitig den Wert für mich erhöht, indem ich das hochgewichtige Objekt C03 (Wichtigkeit 668) einbeziehe, während ich weniger wertvolle Gegenstände für mich weglasse.'\}} \\
            }
        }
    }
    & \\ \\

    \theutterance \stepcounter{utterance}  
    & \multicolumn{4}{p{0.6\linewidth}}{
        \cellcolor[rgb]{0.9,0.9,0.9}{
            \makecell[{{p{\linewidth}}}]{
                \texttt{\tiny{[P1$\rangle$GM]}}
                \texttt{VORSCHLAG: \{'C32', 'A03', 'C00', 'C76', 'C10', 'C81', 'C03', 'C62', 'C98', 'B38', 'A83'\}} \\
\texttt{ARGUMENT: \{'Ich habe den niedrigen Nutzen von C56 (Aufwand 588, Wichtigkeit 65) aus eurem Vorschlag entfernt und dafür mehrere hochgewichtige Gegenstände (C62, C98, B38, A83) hinzugefügt. Der neue Gesamteffort beträgt 2229 ≤ 2307, während die kombinierte Wichtigkeit für mich von 1596 auf 3668 steigt – ein signifikanter Gewinn für beide Seiten, da ihr immer noch fast alle eurer ursprünglichen Gegenstände behaltet.'\}} \\
            }
        }
    }
    & & \\ \\

    \theutterance \stepcounter{utterance}  
    & & \multicolumn{4}{p{0.6\linewidth}}{
        \cellcolor[rgb]{0.9,0.9,0.9}{
            \makecell[{{p{\linewidth}}}]{
                \texttt{\tiny{[GM$\rangle$P2]}}
                \texttt{VORSCHLAG: \{'C32', 'A03', 'C00', 'C76', 'C10', 'C81', 'C03', 'C62', 'C98', 'B38', 'A83'\}} \\
\texttt{ARGUMENT: \{'Ich habe den niedrigen Nutzen von C56 (Aufwand 588, Wichtigkeit 65) aus eurem Vorschlag entfernt und dafür mehrere hochgewichtige Gegenstände (C62, C98, B38, A83) hinzugefügt. Der neue Gesamteffort beträgt 2229 ≤ 2307, während die kombinierte Wichtigkeit für mich von 1596 auf 3668 steigt – ein signifikanter Gewinn für beide Seiten, da ihr immer noch fast alle eurer ursprünglichen Gegenstände behaltet.'\}} \\
            }
        }
    }
    & \\ \\

    \theutterance \stepcounter{utterance}  
    & & & \multicolumn{4}{p{0.6\linewidth}}{
        \cellcolor[rgb]{0.9,0.9,0.9}{
            \makecell[{{p{\linewidth}}}]{
                \texttt{\tiny{[GM$\langle$P2]}}
                \texttt{VORSCHLAG: \{'C32', 'A03', 'C76', 'C10', 'C81', 'C03', 'C98', 'B38', 'A83', 'A07'\}} \\
\texttt{ARGUMENT: \{'Ich habe den Gesamteffort auf 2171 reduziert (unter 2307) und dafür den Gegenstand A07 (Wichtigkeit 500) aufgenommen, während ich weniger für mich wertvolle Gegenstände (C00, C62) weggelassen habe. Damit steigt meine kombinierte Wichtigkeit von 3635 auf 3716, wodurch das Angebot für mich deutlich attraktiver wird, ohne das Aufwandslimit zu überschreiten.'\}} \\
            }
        }
    }
     \\ \\

    \theutterance \stepcounter{utterance}  
    & & \multicolumn{4}{p{0.6\linewidth}}{
        \cellcolor[rgb]{0.9,0.9,0.9}{
            \makecell[{{p{\linewidth}}}]{
                \texttt{\tiny{[P1$\langle$GM]}}
                \texttt{VORSCHLAG: \{'C32', 'A03', 'C76', 'C10', 'C81', 'C03', 'C98', 'B38', 'A83', 'A07'\}} \\
\texttt{ARGUMENT: \{'Ich habe den Gesamteffort auf 2171 reduziert (unter 2307) und dafür den Gegenstand A07 (Wichtigkeit 500) aufgenommen, während ich weniger für mich wertvolle Gegenstände (C00, C62) weggelassen habe. Damit steigt meine kombinierte Wichtigkeit von 3635 auf 3716, wodurch das Angebot für mich deutlich attraktiver wird, ohne das Aufwandslimit zu überschreiten.'\}} \\
            }
        }
    }
    & \\ \\

    \theutterance \stepcounter{utterance}  
    & \multicolumn{4}{p{0.6\linewidth}}{
        \cellcolor[rgb]{0.9,0.9,0.9}{
            \makecell[{{p{\linewidth}}}]{
                \texttt{\tiny{[P1$\rangle$GM]}}
                \texttt{VORSCHLAG: \{'C32', 'A03', 'C10', 'C81', 'C98', 'B38', 'A83', 'A07', 'C62', 'C00', 'C76'\}} \\
\texttt{ARGUMENT: \{'Durch das Entfernen von C03 (geringe Wichtigkeit) und das Hinzufügen von C62 (Wichtigkeit 668), C00 (Wichtigkeit 389) und C76 (Wichtigkeit 138) steigt meine kombinierte Wichtigkeit auf 3759, während der Gesamteffort nur 2042 bleibt – deutlich unter dem Limit von 2307.'\}} \\
            }
        }
    }
    & & \\ \\

    \theutterance \stepcounter{utterance}  
    & & \multicolumn{4}{p{0.6\linewidth}}{
        \cellcolor[rgb]{0.9,0.9,0.9}{
            \makecell[{{p{\linewidth}}}]{
                \texttt{\tiny{[GM$\rangle$P2]}}
                \texttt{VORSCHLAG: \{'C32', 'A03', 'C10', 'C81', 'C98', 'B38', 'A83', 'A07', 'C62', 'C00', 'C76'\}} \\
\texttt{ARGUMENT: \{'Durch das Entfernen von C03 (geringe Wichtigkeit) und das Hinzufügen von C62 (Wichtigkeit 668), C00 (Wichtigkeit 389) und C76 (Wichtigkeit 138) steigt meine kombinierte Wichtigkeit auf 3759, während der Gesamteffort nur 2042 bleibt – deutlich unter dem Limit von 2307.'\}} \\
            }
        }
    }
    & \\ \\

    \theutterance \stepcounter{utterance}  
    & & & \multicolumn{4}{p{0.6\linewidth}}{
        \cellcolor[rgb]{0.9,0.9,0.9}{
            \makecell[{{p{\linewidth}}}]{
                \texttt{\tiny{[GM$\langle$P2]}}
                \texttt{VORSCHLAG: \{'C32', 'A03', 'C76', 'C10', 'C81', 'C03', 'C98', 'B38', 'A83', 'A07'\}} \\
\texttt{ARGUMENT: \{'Dieses Set nutzt exakt 2171 Aufwand (unter 2307) und enthält das für mich höchstwertige Element C03 (Wichtigkeit 668). Im Vergleich zu Ihrem letzten Angebot erhöht sich meine persönliche Wichtigkeit von 3467 auf 3716, während Ihr Nutzen kaum leidet (Sie verlieren C00 mit 389 und C62 mit 30, dafür erhalten Sie den hochgewichtigen C03). Das Ergebnis ist für beide Seiten fast egalwertig, aber insgesamt näher an einer Pareto‑optimalen Verteilung.'\}} \\
            }
        }
    }
     \\ \\

    \theutterance \stepcounter{utterance}  
    & & \multicolumn{4}{p{0.6\linewidth}}{
        \cellcolor[rgb]{0.9,0.9,0.9}{
            \makecell[{{p{\linewidth}}}]{
                \texttt{\tiny{[P1$\langle$GM]}}
                \texttt{VORSCHLAG: \{'C32', 'A03', 'C76', 'C10', 'C81', 'C03', 'C98', 'B38', 'A83', 'A07'\}} \\
\texttt{ARGUMENT: \{'Dieses Set nutzt exakt 2171 Aufwand (unter 2307) und enthält das für mich höchstwertige Element C03 (Wichtigkeit 668). Im Vergleich zu Ihrem letzten Angebot erhöht sich meine persönliche Wichtigkeit von 3467 auf 3716, während Ihr Nutzen kaum leidet (Sie verlieren C00 mit 389 und C62 mit 30, dafür erhalten Sie den hochgewichtigen C03). Das Ergebnis ist für beide Seiten fast egalwertig, aber insgesamt näher an einer Pareto‑optimalen Verteilung.'\}} \\
            }
        }
    }
    & \\ \\

    \theutterance \stepcounter{utterance}  
    & \multicolumn{4}{p{0.6\linewidth}}{
        \cellcolor[rgb]{0.9,0.9,0.9}{
            \makecell[{{p{\linewidth}}}]{
                \texttt{\tiny{[P1$\rangle$GM]}}
                \texttt{VORSCHLAG: \{'C32', 'A03', 'C76', 'C10', 'C81', 'C98', 'B38', 'A83', 'A07', 'C62', 'C00'\}} \\
\texttt{ARGUMENT: \{'Ich ersetze das für mich kaum wertvolle C03 (Wichtigkeit 30, Aufwand 543) durch C62 (Wichtigkeit 668, Aufwand 106) und C00 (Wichtigkeit 389, Aufwand 308). Das erhöht meine kombinierte Wichtigkeit von 2632 auf 3639, während der Gesamteffort nur 2042 beträgt – deutlich unter dem Limit von 2307. Damit erhalte ich deutlich mehr Nutzen, ohne Ihnen wesentliche Gegenstände zu entziehen.'\}} \\
            }
        }
    }
    & & \\ \\

    \theutterance \stepcounter{utterance}  
    & & \multicolumn{4}{p{0.6\linewidth}}{
        \cellcolor[rgb]{0.9,0.9,0.9}{
            \makecell[{{p{\linewidth}}}]{
                \texttt{\tiny{[GM$\rangle$P2]}}
                \texttt{VORSCHLAG: \{'C32', 'A03', 'C76', 'C10', 'C81', 'C98', 'B38', 'A83', 'A07', 'C62', 'C00'\}} \\
\texttt{ARGUMENT: \{'Ich ersetze das für mich kaum wertvolle C03 (Wichtigkeit 30, Aufwand 543) durch C62 (Wichtigkeit 668, Aufwand 106) und C00 (Wichtigkeit 389, Aufwand 308). Das erhöht meine kombinierte Wichtigkeit von 2632 auf 3639, während der Gesamteffort nur 2042 beträgt – deutlich unter dem Limit von 2307. Damit erhalte ich deutlich mehr Nutzen, ohne Ihnen wesentliche Gegenstände zu entziehen.'\}} \\
            }
        }
    }
    & \\ \\

    \theutterance \stepcounter{utterance}  
    & & & \multicolumn{4}{p{0.6\linewidth}}{
        \cellcolor[rgb]{0.9,0.9,0.9}{
            \makecell[{{p{\linewidth}}}]{
                \texttt{\tiny{[GM$\langle$P2]}}
                \texttt{VORSCHLAG: \{'C32', 'A03', 'C76', 'C10', 'C81', 'C98', 'A83', 'A07', 'C00', 'C03'\}} \\
\texttt{ARGUMENT: \{'Dieses Set bleibt unter dem Maximalaufwand von 2307 (2203) und erhöht meine persönliche Wichtigkeit deutlich auf 3984, da ich das hochgewichtige Element C03 (668) hinzufügen konnte, während ich nur wenig wichtige Gegenstände (C62 und B38) entferne, die für mich kaum Nutzen bringen.'\}} \\
            }
        }
    }
     \\ \\

    \theutterance \stepcounter{utterance}  
    & & \multicolumn{4}{p{0.6\linewidth}}{
        \cellcolor[rgb]{0.9,0.9,0.9}{
            \makecell[{{p{\linewidth}}}]{
                \texttt{\tiny{[P1$\langle$GM]}}
                \texttt{VORSCHLAG: \{'C32', 'A03', 'C76', 'C10', 'C81', 'C98', 'A83', 'A07', 'C00', 'C03'\}} \\
\texttt{ARGUMENT: \{'Dieses Set bleibt unter dem Maximalaufwand von 2307 (2203) und erhöht meine persönliche Wichtigkeit deutlich auf 3984, da ich das hochgewichtige Element C03 (668) hinzufügen konnte, während ich nur wenig wichtige Gegenstände (C62 und B38) entferne, die für mich kaum Nutzen bringen.'\}} \\
            }
        }
    }
    & \\ \\

    \theutterance \stepcounter{utterance}  
    & \multicolumn{4}{p{0.6\linewidth}}{
        \cellcolor[rgb]{0.9,0.9,0.9}{
            \makecell[{{p{\linewidth}}}]{
                \texttt{\tiny{[P1$\rangle$GM]}}
                \texttt{VORSCHLAG: \{'C32', 'A03', 'C76', 'C10', 'C81', 'C98', 'A83', 'A07', 'C00', 'C62', 'B38'\}} \\
\texttt{ARGUMENT: \{'Ich ersetze das für mich wenigwertige C03 (Wichtigkeit 30, Aufwand 543) durch C62 (Wichtigkeit 668, Aufwand 106) und füge B38 (Wichtigkeit 500, Aufwand 276) hinzu. Der Gesamteffort beträgt nur 2042 ≤ 2307, während meine kombinierte Wichtigkeit von 2621 auf 3759 steigt. Sie behalten alle anderen Gegenstände, die Sie bereits erhalten haben, und wir nähern uns einer Pareto‑optimalen Verteilung.'\}} \\
            }
        }
    }
    & & \\ \\

    \theutterance \stepcounter{utterance}  
    & & \multicolumn{4}{p{0.6\linewidth}}{
        \cellcolor[rgb]{0.9,0.9,0.9}{
            \makecell[{{p{\linewidth}}}]{
                \texttt{\tiny{[GM$\rangle$P2]}}
                \texttt{VORSCHLAG: \{'C32', 'A03', 'C76', 'C10', 'C81', 'C98', 'A83', 'A07', 'C00', 'C62', 'B38'\}} \\
\texttt{ARGUMENT: \{'Ich ersetze das für mich wenigwertige C03 (Wichtigkeit 30, Aufwand 543) durch C62 (Wichtigkeit 668, Aufwand 106) und füge B38 (Wichtigkeit 500, Aufwand 276) hinzu. Der Gesamteffort beträgt nur 2042 ≤ 2307, während meine kombinierte Wichtigkeit von 2621 auf 3759 steigt. Sie behalten alle anderen Gegenstände, die Sie bereits erhalten haben, und wir nähern uns einer Pareto‑optimalen Verteilung.'\}} \\
            }
        }
    }
    & \\ \\

    \theutterance \stepcounter{utterance}  
    & & & \multicolumn{4}{p{0.6\linewidth}}{
        \cellcolor[rgb]{0.9,0.9,0.9}{
            \makecell[{{p{\linewidth}}}]{
                \texttt{\tiny{[GM$\langle$P2]}}
                \texttt{VORSCHLAG: \{'C32', 'A03', 'C76', 'C10', 'C81', 'C98', 'A83', 'A07', 'C00', 'C03'\}} \\
\texttt{ARGUMENT: \{'Ich schlage vor, dass wir C62 (Wichtigkeit 30, Aufwand 106) und B38 (Wichtigkeit 121, Aufwand 276) aus Ihrem Angebot entfernen und dafür C03 (Wichtigkeit 668, Aufwand 543) aufnehmen. Der Gesamteffort beträgt damit 2203 ≤ 2307. Für mich steigt die kombinierte Wichtigkeit von 3467 auf 3984, während Sie nur Gegenstände mit relativ geringem Nutzen für Sie verlieren.'\}} \\
            }
        }
    }
     \\ \\

    \theutterance \stepcounter{utterance}  
    & & \multicolumn{4}{p{0.6\linewidth}}{
        \cellcolor[rgb]{0.9,0.9,0.9}{
            \makecell[{{p{\linewidth}}}]{
                \texttt{\tiny{[P1$\langle$GM]}}
                \texttt{VORSCHLAG: \{'C32', 'A03', 'C76', 'C10', 'C81', 'C98', 'A83', 'A07', 'C00', 'C03'\}} \\
\texttt{ARGUMENT: \{'Ich schlage vor, dass wir C62 (Wichtigkeit 30, Aufwand 106) und B38 (Wichtigkeit 121, Aufwand 276) aus Ihrem Angebot entfernen und dafür C03 (Wichtigkeit 668, Aufwand 543) aufnehmen. Der Gesamteffort beträgt damit 2203 ≤ 2307. Für mich steigt die kombinierte Wichtigkeit von 3467 auf 3984, während Sie nur Gegenstände mit relativ geringem Nutzen für Sie verlieren.'\}} \\
            }
        }
    }
    & \\ \\

    \theutterance \stepcounter{utterance}  
    & \multicolumn{4}{p{0.6\linewidth}}{
        \cellcolor[rgb]{0.9,0.9,0.9}{
            \makecell[{{p{\linewidth}}}]{
                \texttt{\tiny{[P1$\rangle$GM]}}
                \texttt{VORSCHLAG: \{'C32', 'A03', 'C76', 'C81', 'C98', 'A83', 'A07', 'C00', 'C62', 'B38', 'C03'\}} \\
\texttt{ARGUMENT: \{'Ich biete Ihnen das für Sie hochgewichtige Element C03 (Wichtigkeit 668) an, während ich gleichzeitig C62 (Wichtigkeit 668) und B38 (Wichtigkeit 500) behalte, die für mich entscheidend sind. Der Gesamteffort beträgt 2260 ≤ 2307. Damit erhalten Sie Ihren Top‑Gegenstand, ich erhalte meine beiden wichtigsten Werte, und beide Seiten bleiben unter dem Aufwandslimit.'\}} \\
            }
        }
    }
    & & \\ \\

    \theutterance \stepcounter{utterance}  
    & & \multicolumn{4}{p{0.6\linewidth}}{
        \cellcolor[rgb]{0.9,0.9,0.9}{
            \makecell[{{p{\linewidth}}}]{
                \texttt{\tiny{[GM$\rangle$P2]}}
                \texttt{VORSCHLAG: \{'C32', 'A03', 'C76', 'C81', 'C98', 'A83', 'A07', 'C00', 'C62', 'B38', 'C03'\}} \\
\texttt{ARGUMENT: \{'Ich biete Ihnen das für Sie hochgewichtige Element C03 (Wichtigkeit 668) an, während ich gleichzeitig C62 (Wichtigkeit 668) und B38 (Wichtigkeit 500) behalte, die für mich entscheidend sind. Der Gesamteffort beträgt 2260 ≤ 2307. Damit erhalten Sie Ihren Top‑Gegenstand, ich erhalte meine beiden wichtigsten Werte, und beide Seiten bleiben unter dem Aufwandslimit.'\}} \\
            }
        }
    }
    & \\ \\

    \theutterance \stepcounter{utterance}  
    & & & \multicolumn{4}{p{0.6\linewidth}}{
        \cellcolor[rgb]{0.9,0.9,0.9}{
            \makecell[{{p{\linewidth}}}]{
                \texttt{\tiny{[GM$\langle$P2]}}
                \texttt{VORSCHLAG: \{'C32', 'A03', 'C76', 'C81', 'C98', 'A83', 'A07', 'C03', 'C56'\}} \\
\texttt{ARGUMENT: \{'Ich ersetze die für mich wenigwertigen Gegenstände C62 (Wichtigkeit 30, Aufwand 106) und B38 (Wichtigkeit 121, Aufwand 276) sowie C00 (Wichtigkeit 389, Aufwand 308) durch das hochgewichtige Element C56 (Wichtigkeit 583, Aufwand 588). Der Gesamteffort beträgt jetzt 2158 ≤ 2307, und meine kombinierte Wichtigkeit steigt von 3674 auf 3717, während Sie fast alle Ihrer gewünschten Gegenstände behalten.'\}} \\
            }
        }
    }
     \\ \\

    \theutterance \stepcounter{utterance}  
    & & \multicolumn{4}{p{0.6\linewidth}}{
        \cellcolor[rgb]{0.9,0.9,0.9}{
            \makecell[{{p{\linewidth}}}]{
                \texttt{\tiny{[P1$\langle$GM]}}
                \texttt{VORSCHLAG: \{'C32', 'A03', 'C76', 'C81', 'C98', 'A83', 'A07', 'C03', 'C56'\}} \\
\texttt{ARGUMENT: \{'Ich ersetze die für mich wenigwertigen Gegenstände C62 (Wichtigkeit 30, Aufwand 106) und B38 (Wichtigkeit 121, Aufwand 276) sowie C00 (Wichtigkeit 389, Aufwand 308) durch das hochgewichtige Element C56 (Wichtigkeit 583, Aufwand 588). Der Gesamteffort beträgt jetzt 2158 ≤ 2307, und meine kombinierte Wichtigkeit steigt von 3674 auf 3717, während Sie fast alle Ihrer gewünschten Gegenstände behalten.'\}} \\
            }
        }
    }
    & \\ \\

    \theutterance \stepcounter{utterance}  
    & \multicolumn{4}{p{0.6\linewidth}}{
        \cellcolor[rgb]{0.9,0.9,0.9}{
            \makecell[{{p{\linewidth}}}]{
                \texttt{\tiny{[P1$\rangle$GM]}}
                \texttt{VORSCHLAG: \{'C32', 'A03', 'C76', 'C81', 'C98', 'A83', 'C62', 'B38', 'C00', 'C10'\}} \\
\texttt{ARGUMENT: \{'Ich ersetze die für mich wenig wertvollen Gegenstände C03, C56 und A07 (hoher Aufwand, niedrige Wichtigkeit) durch C62 (Wichtigkeit 668, Aufwand 106), B38 (Wichtigkeit 500, Aufwand 276), C00 (Wichtigkeit 389, Aufwand 308) und C10 (Wichtigkeit 215, Aufwand 325). Der Gesamteffort beträgt nur 1686 ≤ 2307, und meine kombinierte Wichtigkeit steigt von 2082 auf 3638 – ein deutliches Plus für mich, während die von Ihnen gewünschten Gegenstände größtenteils erhalten bleiben.'\}} \\
            }
        }
    }
    & & \\ \\

    \theutterance \stepcounter{utterance}  
    & & \multicolumn{4}{p{0.6\linewidth}}{
        \cellcolor[rgb]{0.9,0.9,0.9}{
            \makecell[{{p{\linewidth}}}]{
                \texttt{\tiny{[GM$\rangle$P2]}}
                \texttt{VORSCHLAG: \{'C32', 'A03', 'C76', 'C81', 'C98', 'A83', 'C62', 'B38', 'C00', 'C10'\}} \\
\texttt{ARGUMENT: \{'Ich ersetze die für mich wenig wertvollen Gegenstände C03, C56 und A07 (hoher Aufwand, niedrige Wichtigkeit) durch C62 (Wichtigkeit 668, Aufwand 106), B38 (Wichtigkeit 500, Aufwand 276), C00 (Wichtigkeit 389, Aufwand 308) und C10 (Wichtigkeit 215, Aufwand 325). Der Gesamteffort beträgt nur 1686 ≤ 2307, und meine kombinierte Wichtigkeit steigt von 2082 auf 3638 – ein deutliches Plus für mich, während die von Ihnen gewünschten Gegenstände größtenteils erhalten bleiben.'\}} \\
            }
        }
    }
    & \\ \\

    \theutterance \stepcounter{utterance}  
    & & & \multicolumn{4}{p{0.6\linewidth}}{
        \cellcolor[rgb]{0.9,0.9,0.9}{
            \makecell[{{p{\linewidth}}}]{
                \texttt{\tiny{[GM$\langle$P2]}}
                \texttt{VORSCHLAG: \{'C32', 'A03', 'C76', 'C81', 'C98', 'A83', 'C10', 'C03', 'B38', 'C00'\}} \\
\texttt{ARGUMENT: \{'Ich ersetze Ihren low‑value Gegenstand C62 (Wichtigkeit 30) durch mein höchstwertiges Element C03 (Wichtigkeit 668). Gleichzeitig behalte ich B38 (Wichtigkeit 121) und C00 (Wichtigkeit 389) bei. Der Gesamteffort beträgt damit 2123 ≤ 2307, und meine kombinierte Wichtigkeit steigt auf 3605, während Sie fast alle Ihrer gewünschten Objekte erhalten und zusätzlich den Top‑Gegenstand C03 von mir bekommen.'\}} \\
            }
        }
    }
     \\ \\

    \theutterance \stepcounter{utterance}  
    & & \multicolumn{4}{p{0.6\linewidth}}{
        \cellcolor[rgb]{0.9,0.9,0.9}{
            \makecell[{{p{\linewidth}}}]{
                \texttt{\tiny{[P1$\langle$GM]}}
                \texttt{VORSCHLAG: \{'C32', 'A03', 'C76', 'C81', 'C98', 'A83', 'C10', 'C03', 'B38', 'C00'\}} \\
\texttt{ARGUMENT: \{'Ich ersetze Ihren low‑value Gegenstand C62 (Wichtigkeit 30) durch mein höchstwertiges Element C03 (Wichtigkeit 668). Gleichzeitig behalte ich B38 (Wichtigkeit 121) und C00 (Wichtigkeit 389) bei. Der Gesamteffort beträgt damit 2123 ≤ 2307, und meine kombinierte Wichtigkeit steigt auf 3605, während Sie fast alle Ihrer gewünschten Objekte erhalten und zusätzlich den Top‑Gegenstand C03 von mir bekommen.'\}} \\
            }
        }
    }
    & \\ \\

    \theutterance \stepcounter{utterance}  
    & \multicolumn{4}{p{0.6\linewidth}}{
        \cellcolor[rgb]{0.9,0.9,0.9}{
            \makecell[{{p{\linewidth}}}]{
                \texttt{\tiny{[P1$\rangle$GM]}}
                \texttt{VORSCHLAG: \{'C32', 'A03', 'C76', 'C81', 'C98', 'A83', 'C10', 'C03', 'B38', 'C00', 'C62'\}} \\
\texttt{ARGUMENT: \{'Ich behalte alle von Ihnen bereits vorgeschlagenen Gegenstände, die für mich einen hohen Nutzen bringen, und füge zusätzlich C62 (Wichtigkeit 668, Aufwand 106) hinzu, das für mich den größten Wert liefert. Der Gesamteffort beträgt 2229 ≤ 2307, sodass das Limit eingehalten wird. Sie verlieren lediglich C03 (Wichtigkeit 30 für Sie), während Sie weiterhin fast alle Ihrer gewünschten Gegenstände erhalten.'\}} \\
            }
        }
    }
    & & \\ \\

    \theutterance \stepcounter{utterance}  
    & & \multicolumn{4}{p{0.6\linewidth}}{
        \cellcolor[rgb]{0.9,0.9,0.9}{
            \makecell[{{p{\linewidth}}}]{
                \texttt{\tiny{[GM$\rangle$P2]}}
                \texttt{VORSCHLAG: \{'C32', 'A03', 'C76', 'C81', 'C98', 'A83', 'C10', 'C03', 'B38', 'C00', 'C62'\}} \\
\texttt{ARGUMENT: \{'Ich behalte alle von Ihnen bereits vorgeschlagenen Gegenstände, die für mich einen hohen Nutzen bringen, und füge zusätzlich C62 (Wichtigkeit 668, Aufwand 106) hinzu, das für mich den größten Wert liefert. Der Gesamteffort beträgt 2229 ≤ 2307, sodass das Limit eingehalten wird. Sie verlieren lediglich C03 (Wichtigkeit 30 für Sie), während Sie weiterhin fast alle Ihrer gewünschten Gegenstände erhalten.'\}} \\
            }
        }
    }
    & \\ \\

    \theutterance \stepcounter{utterance}  
    & & & \multicolumn{4}{p{0.6\linewidth}}{
        \cellcolor[rgb]{0.9,0.9,0.9}{
            \makecell[{{p{\linewidth}}}]{
                \texttt{\tiny{[GM$\langle$P2]}}
                \texttt{VORSCHLAG: \{'C32', 'A03', 'C76', 'C81', 'C98', 'A83', 'C10', 'C03', 'C56'\}} \\
\texttt{ARGUMENT: \{'Ich ersetze die für mich wenigwertigen Gegenstände C62 (Wichtigkeit 30, Aufwand 106), C00 (Wichtigkeit 389, Aufwand 308) und B38 (Wichtigkeit 121, Aufwand 276) durch das hochgewichtige Element C56 (Wichtigkeit 583, Aufwand 588). Damit bleibt der Gesamteffort bei 2127 ≤ 2307 und meine kombinierte Wichtigkeit steigt von 3635 auf 3678, während Sie fast alle Ihrer gewünschten Gegenstände behalten.'\}} \\
            }
        }
    }
     \\ \\

    \theutterance \stepcounter{utterance}  
    & & \multicolumn{4}{p{0.6\linewidth}}{
        \cellcolor[rgb]{0.9,0.9,0.9}{
            \makecell[{{p{\linewidth}}}]{
                \texttt{\tiny{[P1$\langle$GM]}}
                \texttt{VORSCHLAG: \{'C32', 'A03', 'C76', 'C81', 'C98', 'A83', 'C10', 'C03', 'C56'\}} \\
\texttt{ARGUMENT: \{'Ich ersetze die für mich wenigwertigen Gegenstände C62 (Wichtigkeit 30, Aufwand 106), C00 (Wichtigkeit 389, Aufwand 308) und B38 (Wichtigkeit 121, Aufwand 276) durch das hochgewichtige Element C56 (Wichtigkeit 583, Aufwand 588). Damit bleibt der Gesamteffort bei 2127 ≤ 2307 und meine kombinierte Wichtigkeit steigt von 3635 auf 3678, während Sie fast alle Ihrer gewünschten Gegenstände behalten.'\}} \\
            }
        }
    }
    & \\ \\

    \theutterance \stepcounter{utterance}  
    & \multicolumn{4}{p{0.6\linewidth}}{
        \cellcolor[rgb]{0.9,0.9,0.9}{
            \makecell[{{p{\linewidth}}}]{
                \texttt{\tiny{[P1$\rangle$GM]}}
                \texttt{VORSCHLAG: \{'C03', 'C62', 'B38', 'C00', 'C98', 'A83', 'C32', 'A03', 'C76', 'C10', 'C81'\}} \\
\texttt{ARGUMENT: \{'Dieses Set erfüllt das Aufwandslimit (2229 ≤ 2307) und enthält Ihr wichtiges Element C03, das Sie sehr hoch gewichten. Gleichzeitig behalte ich meine wertvollsten Gegenstände (C62, B38, C00, C98, A83, C32, A03, C76, C10, C81), sodass meine kombinierte Wichtigkeit auf 3691 steigt. Damit erhalten beide Seiten ihre Kernobjekte und wir bleiben deutlich unter dem Maximalaufwand.'\}} \\
            }
        }
    }
    & & \\ \\

    \theutterance \stepcounter{utterance}  
    & & \multicolumn{4}{p{0.6\linewidth}}{
        \cellcolor[rgb]{0.9,0.9,0.9}{
            \makecell[{{p{\linewidth}}}]{
                \texttt{\tiny{[GM$\rangle$P2]}}
                \texttt{VORSCHLAG: \{'C03', 'C62', 'B38', 'C00', 'C98', 'A83', 'C32', 'A03', 'C76', 'C10', 'C81'\}} \\
\texttt{ARGUMENT: \{'Dieses Set erfüllt das Aufwandslimit (2229 ≤ 2307) und enthält Ihr wichtiges Element C03, das Sie sehr hoch gewichten. Gleichzeitig behalte ich meine wertvollsten Gegenstände (C62, B38, C00, C98, A83, C32, A03, C76, C10, C81), sodass meine kombinierte Wichtigkeit auf 3691 steigt. Damit erhalten beide Seiten ihre Kernobjekte und wir bleiben deutlich unter dem Maximalaufwand.'\}} \\
            }
        }
    }
    & \\ \\

    \theutterance \stepcounter{utterance}  
    & & & \multicolumn{4}{p{0.6\linewidth}}{
        \cellcolor[rgb]{0.9,0.9,0.9}{
            \makecell[{{p{\linewidth}}}]{
                \texttt{\tiny{[GM$\langle$P2]}}
                \texttt{ZUSTIMMUNG: \{'C03', 'C62', 'B38', 'C00', 'C98', 'A83', 'C32', 'A03', 'C76', 'C10', 'C81'\}} \\
\texttt{ARGUMENT: \{'Ich akzeptiere Ihr Angebot, da es unter dem Aufwandslimit liegt und mir einen hohen Gesamtnutzen von 3635 liefert, während Sie Ihre wichtigsten Gegenstände behalten.'\}} \\
            }
        }
    }
     \\ \\

    \theutterance \stepcounter{utterance}  
    & & & \multicolumn{2}{p{0.3\linewidth}}{
        \cellcolor[rgb]{0.9,0.9,0.9}{
            \makecell[{{p{\linewidth}}}]{
                \texttt{\tiny{[GM$|$GM]}}
                \texttt{game successful} \\
            }
        }
    }
    & & \\ \\

    \theutterance \stepcounter{utterance}  
    & & & \multicolumn{2}{p{0.3\linewidth}}{
        \cellcolor[rgb]{0.9,0.9,0.9}{
            \makecell[{{p{\linewidth}}}]{
                \texttt{\tiny{[GM$|$GM]}}
                \texttt{end game} \\
            }
        }
    }
    & & \\ \\

\end{supertabular}
}

\end{document}
